\newcommand{\Tone}{\begin{pmatrix}1&2\\0&1\end{pmatrix}}
\newcommand{\Ttwo}{\begin{pmatrix}1&1\\2&0\end{pmatrix}}
\newcommand{\Tthree}{\begin{pmatrix}0&1\\1&1\end{pmatrix}}
\newcommand{\Tfour}{\begin{pmatrix}2&1\\1&0\end{pmatrix}}
\newcommand{\Mtwelve}{\begin{pmatrix}\colorbox{tngreen!22}{$5$}&\colorbox{tngreen!22}{$1$}\\\colorbox{tngreen!22}{$2$}&\colorbox{tngreen!22}{$0$}\end{pmatrix}}
\newcommand{\Mthirtyfour}{\begin{pmatrix}\colorbox{tngreen!22}{$1$}&\colorbox{tngreen!22}{$0$}\\\colorbox{tngreen!22}{$3$}&\colorbox{tngreen!22}{$1$}\end{pmatrix}}
\newcommand{\Zvalue}{8}
\newcommand{\Kthirteen}{\begin{pmatrix}0&\colorbox{tnorange!22}{$1$}&\colorbox{tnorange!22}{$1$}&\colorbox{tnorange!22}{$1$}\\0&\colorbox{tnorange!22}{$2$}&\colorbox{tnorange!22}{$2$}&\colorbox{tnorange!22}{$2$}\\0&0&0&0\\0&\colorbox{tnorange!22}{$1$}&\colorbox{tnorange!22}{$1$}&\colorbox{tnorange!22}{$1$}\end{pmatrix}}
\newcommand{\Ktwentyfour}{\begin{pmatrix}\colorbox{tnorange!22}{$2$}&\colorbox{tnorange!22}{$1$}&\colorbox{tnorange!22}{$2$}&\colorbox{tnorange!22}{$1$}\\\colorbox{tnorange!22}{$4$}&\colorbox{tnorange!22}{$2$}&0&0\\\colorbox{tnorange!22}{$1$}&0&\colorbox{tnorange!22}{$1$}&0\\\colorbox{tnorange!22}{$2$}&0&0&0\end{pmatrix}}
\newcommand{\Kentries}{9/9}
\newcommand{\Mtwelveexample}{1\cdot1 + 2\cdot2 = 5}
\newcommand{\Cqzero}{\begin{pmatrix}0&0&0&0\\\colorbox{tnblue!22}{$1$}&0&0&0\end{pmatrix}}
\newcommand{\Cqlam}{\begin{pmatrix}\colorbox{tnblue!22}{$1$}&0&0&\colorbox{tnblue!22}{$1$}\\0&0&0&0\end{pmatrix}}
\newcommand{\CqW}{\begin{pmatrix}0&0&0&0\\0&0&0&\colorbox{tnorange!22}{$-1$}\end{pmatrix}}
\newcommand{\Cqnonzero}{8}
\newcommand{\Cqneg}{2}
\newcommand{\Phimatrix}{\begin{pmatrix}1&1\\1&\colorbox{tnorange!22}{$-1$}\end{pmatrix}}
\newcommand{\toyOA}{725571/800000}
\newcommand{\toyOB}{19701/800000}
\newcommand{\toyPL}{199/7329}
\newcommand{\toyW}{\begin{pmatrix}9/10&1/20\\1/40&1/40\end{pmatrix}}
\newcommand{\toyFA}{\begin{pmatrix}1&1/2\\0&0\end{pmatrix}}
\newcommand{\toyFB}{\begin{pmatrix}0&1/2\\0&0\end{pmatrix}}
\newcommand{\joinentries}{4}

\providecommand{\hl}[2]{\colorbox{#1!22}{#2}}
\section{Contraction orders and the growth average}\label{app:tensor-examples-v2}

This appendix describes two things the later appendices require. First, a contraction order and its tree, shown on a toy example of four contracted tensors. Then the growth average written as a sum over fault configurations, which Appendix~\ref{app:growth} evaluates as one contraction.
\subsection{Two orders for the same contraction}
A contraction order specifies which tensors to combine first. Draw a bubble around a group, multiply their entries, sum the indices used only inside that group, and retain the collective tensor legs. 
This results in one tensor that can then be combined with the next group. 
In Fig.~\ref{fig:contraction-tree} the four indices are binary and the complete scalar is
\begin{equation}
 Z=\sum_{a,b,c,d}T_1(a,b)T_2(b,c)T_3(c,d)T_4(d,a).
 \label{eq:treeexample}
\end{equation}
{\color{black}
\diagram[H]{contraction_tree}
{Two orders for the same network. Each bubble becomes one tensor, and its exposed indices determine its size.}
{fig:contraction-tree}}
\FloatBarrier

For example, take the tensors, with (rows, columns) being the (first, second) indices,
\[
T_1=\Tone,\quad T_2=\Ttwo,\quad T_3=\Tthree,\quad T_4=\Tfour .
\]
Joining neighbours first sums $b$ and $d$,
\begin{equation}
 \begin{aligned}
 M_{12}(a,c)&=\sum_bT_1(a,b)T_2(b,c)=\Mtwelve,\\
 M_{34}(c,a)&=\sum_dT_3(c,d)T_4(d,a)=\Mthirtyfour,\\
 Z&=\sum_{a,c}M_{12}(a,c)M_{34}(c,a)=\Zvalue ,
 \end{aligned}
 \label{eq:treejoins}
\end{equation}
for instance $M_{12}(0,0)=\Mtwelveexample$. Each intermediate has four entries. Joining opposite tensors first sums nothing, because $T_1$ and $T_3$ share no index, so the intermediates keep all four legs open. Each of $K_{13}(a,b,c,d)=T_1(a,b)T_3(c,d)$ and $K_{24}(a,b,c,d)=T_2(b,c)T_4(d,a)$ then has sixteen entries, and $Z=\sum_{a,b,c,d}K_{13}K_{24}=\Zvalue$ again. The order changes the storage and the arithmetic but both calculations sum to the same $Z$.
{\color{black}
\begin{center}
\begin{tikzpicture}[x=1cm,y=1cm]
\node[tn tensor,fill=tnorange,minimum width=8mm,minimum height=8mm] (t1) at (0,1.6) {$T_1$};
\node[tn tensor,fill=tnorange,minimum width=8mm,minimum height=8mm] (t2) at (2.2,1.6) {$T_2$};
\node[tn tensor,fill=tnorange,minimum width=8mm,minimum height=8mm] (t4) at (0,0) {$T_4$};
\node[tn tensor,fill=tnorange,minimum width=8mm,minimum height=8mm] (t3) at (2.2,0) {$T_3$};
\draw[tn leg] (t1)--node[above,tn label] {$b$}(t2); \draw[tn leg] (t4)--node[below,tn label] {$d$}(t3);
\draw[tn leg] (t1)--node[left,tn label] {$a$}(t4); \draw[tn leg] (t2)--node[right,tn label] {$c$}(t3);
\node[font=\Large] at (-1.6,0.8) {$\displaystyle\sum_{a,b,c,d}$};
\node[font=\large] at (3.4,0.8) {$=$};
\node[font=\Large] at (4.4,0.8) {$\displaystyle\sum_{a,c}$};
\node[tn intermediate,minimum width=12mm,minimum height=11mm] (m12) at (5.6,0.8) {$M_{12}$};
\node[tn intermediate,minimum width=12mm,minimum height=11mm] (m34) at (7.6,0.8) {$M_{34}$};
\draw[tn leg] ([yshift=2.5mm]m12.east)--node[above,tn label] {$a$}([yshift=2.5mm]m34.west);
\draw[tn leg] ([yshift=-2.5mm]m12.east)--node[below,tn label] {$c$}([yshift=-2.5mm]m34.west);
\node[font=\large] at (8.7,0.8) {$=\ \Zvalue$};
\end{tikzpicture}
\end{center}}

The order in which tensors are joined does not change the answer but greatly affects the size of the intermediate tables. For a large network, contraction order decides whether the contraction fits in memory at all. 
The growth stage (see Fig.~\ref{fig:overview}) has 569 factors, so the contraction order decides how large the intermediate arrays blow up to. The plan used for the growth sum keeps the largest stored table at $2^{23}$ entries while the largest single join multiplies $2^{27}$ pairs of entries before summing them.

\subsection{Including the noisy growth stage}\label{app:growth-overview-v2}
The earlier $d=3$ stage supplies weights for its outgoing syndrome and logical Pauli label. Clifford growth propagates these labels and adds the labels of its faults. They combine by XOR: two faults may cancel a detector mismatch, so they must not be rejected separately. For a growth-fault configuration $g$, let $w_g$ be its probability, $h(g)$ its detector-compatibility label, $s(g)$ its required input syndrome, and $v(g)$ its output syndrome and logical frame. Write $W_{s,t}$ for the earlier-stage weight with syndrome $s$ and logical Pauli label $t$. Then
\begin{equation}
 \begin{gathered}
 O=\sum_{g,t}w_g\,\ind{h(g)=0}\,
       W_{s(g),t}\,F_O\!\left(\Gamma(v(g),t)\right),\\
 O\in\{A,B\}.
 \end{gathered}
 \label{eq:physicalgrowth}
\end{equation}
The map $\Gamma$ retains the outgoing syndrome and XORs $t$ into the logical frame. This is Eq.~\eqref{eq:responseaverage} with the earlier weights written explicitly. 
Reading Eq.~\eqref{eq:physicalgrowth} from right to left, $F_O$ is the final $d=5$ response table, indexed by the label that arrives at the last check. $W_{s,t}$ is the first-stage table, first-stage post-selection included. 
In between, a configuration $g$ happens with probability $w_g$, is kept only if its detector label is zero, needs the input syndrome $s(g)$, and emits $v(g)$. 

As a toy example, take two syndromes, two tags and four growth configurations at $p=10^{-2}$. Rows are indexed by the syndrome, columns by the tag of $W$ or by the frame of $F_O$, so that
\[
W=\toyW,\qquad F_A=\toyFA,\qquad F_B=\toyFB .
\]
Each configuration contributes one row per tag:
{\color{black}
\begin{center}\small\begin{tabular}{ccccccc}\toprule
$g$ & $t$ & $w_g$ & $h(g)$ & $s(g)$ & $\Gamma(v(g),t)$ & contribution to $A$, to $B$ \\\midrule
1 & 0 & $9801/10000$ & 0 & 0 & $(0,0)$ & $88209/100000$, $0$ \\
1 & 1 & $9801/10000$ & 0 & 0 & $(0,1)$ & $9801/400000$, $9801/400000$ \\
2 & 0 & $99/10000$ & 0 & 1 & $(0,1)$ & $99/800000$, $99/800000$ \\
2 & 1 & $99/10000$ & 0 & 1 & $(0,0)$ & $99/400000$, $0$ \\
3 & 0 & $99/10000$ & 1 & 0 & $(1,0)$ & $0$, $0$ \\
3 & 1 & $99/10000$ & 1 & 0 & $(1,1)$ & $0$, $0$ \\
4 & 0 & $1/10000$ & 0 & 0 & $(1,0)$ & $0$, $0$ \\
4 & 1 & $1/10000$ & 0 & 0 & $(1,1)$ & $0$, $0$ \\
\bottomrule\end{tabular}
\end{center}}
Summing the last column gives $A=\toyOA$ and $B=\toyOB$. Configuration 3 is dropped by $h(g)=1$ whatever its weight, configuration 2 only ever meets the $s=1$ row of $W$, its frame flip moves the tag before it reaches $F_O$, and configuration 4 is compatible but lands on output syndrome 1, where both $F_A$ and $F_B$ vanish in this toy example. 

Two things the formula insists on.
\begin{enumerate}
\item All 64 input syndromes remain in $W$: a growth fault can cancel an input syndrome, as configuration 2 does, so the syndrome-0 row alone would miss accepted events. 
\item Growth faults combine by XOR with each other and with the earlier label before anything is rejected. 
\end{enumerate}
Growth removes cross-syndrome coherence in the way proved in Appendix~\ref{app:growth}, which is what makes a table of probabilities $W_{s,t}$ sufficient. 
Nothing is post-selected between the two stages. Every first-stage syndrome enters growth, and the noiseless final code projection of the standalone $d=3$ calculation is not applied at this interface.
\diagram[H]{growth_contraction}
{The $d=5$ tensor network and one XOR contraction. Rounded boxes are tensors and the arrow denotes a calculation step.}
{fig:growth-network}

Fig.~\ref{fig:growth-network}(a) groups the growth noise into 567 distributions $g_i$ between $W$ and the final response $F_O$. The grey constraint matches their labels by requiring a zero total XOR. Panel (b) combines two partial tables $M_L(u),M_R(v)$ by summing over $u\oplus v=z$ and keeping the output label $z$. With two-bit labels and three entries each,
{\color{black}
\begin{center}\small\begin{tabular}{cccc}\toprule
$z$ & pairs $(u,v)$ with $u\oplus v=z$ & $\sum M_L(u)M_R(v)$ & value \\\midrule
$(0,0)$ & $(00,00)$, $(10,10)$ & $9/10\cdot4/5 + 1/20\cdot1/10$ & $29/40$ \\
$(0,1)$ & $(10,11)$, $(01,00)$ & $1/20\cdot1/10 + 1/20\cdot4/5$ & $9/200$ \\
$(1,0)$ & $(00,10)$, $(10,00)$, $(01,11)$ & $9/10\cdot1/10 + 1/20\cdot4/5 + 1/20\cdot1/10$ & $27/200$ \\
$(1,1)$ & $(00,11)$, $(01,10)$ & $9/10\cdot1/10 + 1/20\cdot1/10$ & $19/200$ \\
\bottomrule\end{tabular}
\end{center}}
the parent table has \joinentries\ entries, one per reachable $z$. Only labels that can be cancelled by the remaining factors are stored. Appendix~\ref{app:growth} gives the binary constraints, contraction and exact Fourier slicing over all characters
